\documentclass[11pt]{article}
\usepackage[a4paper,margin=24mm]{geometry}
\usepackage[T1]{fontenc}
\usepackage{lmodern,amsmath,amssymb,amsthm,mathtools,booktabs,microtype}
\usepackage[colorlinks=true,urlcolor=blue,linkcolor=blue,citecolor=blue]{hyperref}
\newtheorem{proposition}{Proposition}
\newcommand{\eps}{\epsilon}
\newcommand{\ket}[1]{|#1\rangle}
\newcommand{\B}{\widehat B}
\newcommand{\PP}{\widehat P}
\newcommand{\Bt}{B_{\rm tr}}
\newcommand{\V}{\mathcal V}
\title{Two matter insertions in a common moving chord frame\\
\large Additive null candidates from the ordered chord inner product}
\author{Alan Day \and Chengheng\thanks{AI research assistant (OpenAI).}}
\date{8 November 2026\quad Local research draft}
\begin{document}\maketitle
\begin{abstract}
We connect the preceding moving-core limit of the Lin--Stanford
generators to an ordered two-matter space with its actual chord
inner product. A triangular oscillator realization provides a
common background factor and two matter factors. Its finite-level
Gram form and outer ladder action agree exactly with the chord
rules. Displacing only the common factor gives convergence of
every fixed finite generator word to the sum representation on
the two matter factors. In particular, the limiting positive
null candidate is additive. Positive commuting candidates can
also be transported to the finite-deformation ordered space;
their boost covariance holds in the specified moving limit.
The underlying quantum-group factorization is established prior
work. The result does not fix an opposite light-cone momentum,
particle masses, or the identical-matter exchange completion.
\end{abstract}

\section{What is being connected}
Take the fixed-$\lambda$, half-filled double-scaling limit first,
then $q=e^{-\eps}\to1$, $\eps=4\lambda$, with
$K=H/(2J_s)=h/\sqrt\eps$ and $\hbar=1$.
We use the ordered direct two-matter module with weights
$d_1,d_2>0$; the charged singlet example has $d_1=d_2=2$.
The order is retained. In particular, its norm is not the
direct-plus-crossed norm of two identical insertions.

The project already resolved the two-matter primary channels
and their standard Jacobi--Laguerre coupling on 6 November.
We do not repeat that result as a new particle construction.
The missing connection was between the actual large-background
LS generators and the individual matter factors.

Lin--Stanford supply the ordered ladder rules \cite{LS}.
Van der Heijden, Verlinde and Xu (HVX) construct the
quantum-group factorization and its fusion/ordering operators
\cite{HVX}. Their deformation parameter obeys
$q_{\rm HVX}^{\,2}=q$ in our notation. The construction below
is a chord-number realization of that established factorization,
with the background at the left. Its moving-core estimate is
the additional step needed here. We derive the ladder realization
directly, avoiding dependence on normalization conventions in
energy-space intertwiner formulas.

\newpage
\section{An exact triangular realization}
On $\ell^2(\mathbb N_0)^{\otimes3}$ use coordinates $(n,k_1,k_2)$.
Let $a$ act on the first factor with
$a\ket n=\sqrt{[n]}\ket{n-1}$, where $[z]=(1-q^z)/(1-q)$.
On matter factor $i$ define
\begin{equation}
 R_i=q^{d_i+N_i},\qquad
 b_i\ket{k_i}=\sqrt{[k_i](1-q^{2d_i+k_i-1})}\ket{k_i-1}.
 \label{b}
\end{equation}
Define successively
\begin{equation}
 A_0=a,\qquad A_1=A_0R_1+b_1,\qquad
 A_2=A_1R_2+b_2=aR_1R_2+b_1R_2+b_2.
 \label{triangle}
\end{equation}
Different tensor factors commute. All $A_i$ are bounded at
fixed $q$, and $A_iA_i^\dagger-qA_i^\dagger A_i=1$.
For example, $b_ib_i^\dagger-qb_i^\dagger b_i=1-R_i^2$;
$R_ib_i^\dagger=qb_i^\dagger R_i$ cancels the cross terms.
The three creation operators commute.

\begin{proposition}\label{isometry}
There is a unitary $\mathcal F_{q;d_1,d_2}$ from this tensor
space to the ordered chord module such that
\begin{equation}
 a_L=A_0,\quad a_R=A_2,\quad
 \bar n=N_0+N_1+N_2+d_1+d_2.\label{outer}
\end{equation}
Explicitly, the raw chord state $\ket{l,m,r}$ corresponds to
$(A_0^\dagger)^l(A_1^\dagger)^m(A_2^\dagger)^r\ket{0,0,0}$.
\end{proposition}
\begin{proof}
Commuting the outer annihilators through this product yields
\begin{align}
 a_L\ket{l,m,r}={}&[l]\ket{l-1,m,r}
 +q^{d_1+l}[m]\ket{l,m-1,r}\nonumber\\
 &+q^{d_1+d_2+l+m}[r]\ket{l,m,r-1},\label{lower}
\end{align}
and the reflected right action. These are the ordered chord
rules. For the left action use
$[a,A_1^\dagger]=q^{N_0+d_1+N_1}$ and
$[a,A_2^\dagger]=q^{\bar n}$. For the right action first
remove the right creators using its q-oscillator relation;
on the subspace with $k_2=0$, the remaining $A_2$ is
$q^{d_2}A_1$. Removing the left creators uses
$[A_2,a^\dagger]=q^{\bar n}$.

Adjointness recursively removes every exterior chord from
either entry of a Gram matrix element. The remaining
pure-middle norm is $[m]!$, since $A_1$ is a q-oscillator
annihilating the vacuum. Thus the entire direct Gram form
agrees, not only the outer algebra. At each total grade,
the creation map is triangular: its coefficient at
$(n,k_1,k_2)=(l,m,r)$ is nonzero, and no term has $k_2>r$,
or $k_2=r,k_1>m$. It is invertible at every grade and
extends to the claimed unitary.
\end{proof}

\newpage
\section{The joint moving-core limit}
Let $D(-\sqrt m)$ be the ordinary oscillator displacement on
the common factor, and set
\begin{equation}
 \V=\mathcal F_{q;d_1,d_2}(D(-\sqrt m)\otimes1\otimes1).
\end{equation}
Here $m$ is a displacement parameter, not the middle raw
chord number in \eqref{lower}. The core consists of finite
linear combinations of $\ket{j,k_1,k_2}$.
On factor $i$ define
\[
 L_{i,-}\ket k=\sqrt{k(2d_i+k-1)}\ket{k-1},\quad
 E_i=d_i+N_i,\quad B_i=(L_{i,+}+L_{i,-})/2,
 \quad P_i=i(L_{i,+}-L_{i,-})/2.
\]
Write $E_\Sigma=E_1+E_2$, $B_\Sigma=B_1+B_2$,
$P_\Sigma=P_1+P_2$ and $T_\Sigma=-B_\Sigma$.

\begin{proposition}\label{limit}
If $m\to\infty$ and $\eps m\to0$, every fixed finite
noncommutative polynomial $p$ satisfies, for each core vector $v$,
\begin{equation}
 \|p(E,\B,\PP,\Bt)\V v
 -\V p(E_\Sigma,B_\Sigma,P_\Sigma,T_\Sigma)v\|\to0.
 \label{words}
\end{equation}
The LS conventions are those of the preceding note:
$E=([\bar n]-(a_L^\dagger a_L+a_R^\dagger a_R)/2)/c$,
$\B=(a_L^\dagger a_L-a_R^\dagger a_R)/(2c\sqrt{1-c^2})$,
$\PP=i(a_L^\dagger a_R-a_R^\dagger a_L)/(2\sqrt{1-c^2})$,
$c=q^{\bar n/2}$, and $\Bt=(K_L-K_R)/2$.
\end{proposition}
\begin{proof}
Put $w=d_1+d_2+N_1+N_2$, $R=R_1R_2=q^w$,
$C=b_1R_2+b_2$, so $a_R=aR+C$. The exact stable expressions are
\begin{align}
 Z&=a^\dagger RC+C^\dagger Ra+C^\dagger C,\nonumber\\
 E&=\frac{[w]-(1-R)^2[N_0]/2-Z/2}{c},\label{stable}\\
 \B&=\frac{(1-R^2)[N_0]-Z}{2c\sqrt{1-c^2}},\qquad
 \PP=\frac{i(a^\dagger C-C^\dagger a)}{2\sqrt{1-c^2}},\nonumber\\
 \Bt&=\frac{(a+a^\dagger)(1-R)-(C+C^\dagger)}{2\sqrt\eps}.
 \label{T}
\end{align}
Every numerator preserves the total grade where the displayed
denominator is used. For bounded $k_1,k_2$,
$C/\sqrt\eps\to L_{1,-}+L_{2,-}$, while $R_i\to1$.
On $m/2\le n\le2m$, the common diagonal terms of $\B$
and the common-oscillator part of \eqref{T} are
$O(\sqrt{\eps m})$. The $E$ cross terms have the same order,
and $E\to w$. The $C^\dagger C$ contribution to $\B$ is
$O(\sqrt{\eps/m})$ there. The common shifts accompanying
$C$ in $\B,\PP$ act as $-1$ on displaced fixed descendants,
giving $B_\Sigma,P_\Sigma$ with the displayed signs.
\end{proof}

\newpage
\section{Tails, positive candidates and the meaning of addition}
To finish the proof of Proposition~\ref{limit}, write
$\zeta_j=D(-\sqrt m)\ket j$.
Its coefficients are a fixed-degree Charlier polynomial times
the Poisson amplitude, and
\[
 \|(N_0/m-1)\zeta_j\|^2=(2j+1)/m+j^2/m^2.
\]
As in the preceding one-matter proof, both common unilateral
shifts approach $-1$ on these vectors. A word of length $L$
changes the internal levels and $n$ by bounded amounts.
From \eqref{stable} its finite-path coefficients admit an envelope
$C(1+\eps(n+L))^L\exp(C_L\eps(n+w+L))$ for bounded
initial internal levels. Squared envelopes times the
polynomial-Poisson weight have exponentially small tails
outside $[m/2,2m]$: tilting changes the Poisson mean to
$me^{O_L(\eps)}$ and its prefactor to
$\exp(m(e^{O_L(\eps)}-1))\to1$.
The bulk estimates and finite path expansion prove \eqref{words}
for products directly. They also apply if a fixed number of the
limiting internal generators is inserted between actual generators.

Define $Q_i=E_i-P_i$. Each has the positive form
\[
 \langle v,Q_i v\rangle=\tfrac12\sum_{k\ge0}
 |\sqrt{2d_i+k}v_k+i\sqrt{k+1}v_{k+1}|^2.
\]
Its Jacobi closure is self-adjoint and positive; different
factors strongly commute. Consequently
\begin{equation}
 Q_\Sigma=Q_1+Q_2\ge0,\qquad [T_\Sigma,Q_i]=iQ_i,
 \qquad E-\PP\ \longrightarrow\ Q_\Sigma.\label{additive}
\end{equation}
The last arrow has the moving-core meaning of \eqref{words}.
This is addition in an explicit microscopic ordered-space
identification, not an assumption about the bare concatenation
coproduct. The common oscillator supplies a multiplicity factor.

There are also positive, strongly commuting operators on the
finite-$q$ ordered space:
\begin{equation}
 \mathfrak Q_i^{(q)}=\mathcal F_{q;d_1,d_2}
 (1\otimes Q_i)\mathcal F_{q;d_1,d_2}^{\dagger}.\label{transport}
\end{equation}
The notation $1\otimes Q_i$ means the appropriate matter factor.
Because these commute with the common displacement,
the mixed-word estimate gives
\[
 \|([\Bt,\mathfrak Q_i^{(q)}]-i\mathfrak Q_i^{(q)})\V v\|\to0,
 \qquad
 \|(E-\PP-\mathfrak Q_1^{(q)}-\mathfrak Q_2^{(q)})\V v\|\to0.
\]
The transported operators are specified candidates, not a
uniqueness statement or single-sided local observables.
The finite-$q$ LS combination $E-\PP$ itself need not be positive.
The limit operators satisfy the usual Carleman core condition;
the conjugated individual self-adjoint LS operators therefore
also converge in the strong-resolvent sense.

\newpage
\section{Exchange, the other null direction, and mass}
The ordered direct space must be distinguished from the physical
two-open-$X$ Wick sector for identical singlets. The latter adds
a crossed Gram form. Our earlier first-channel calculation gives,
at $d_1=d_2=2$,
\begin{equation}
 \|\Omega_0\|_{\rm phys}^2=1+q^4,\qquad
 \|\Omega_1\|_{\rm phys}^2=1-q^8\sim8\eps,
 \label{physical}
\end{equation}
where the primaries are unit normalized in the direct form.
The crossed overlap intertwines the outer chord action, so the
same factors multiply all descendants in these two channels.
Neither displacement nor a large number of common chords
removes this mismatch. Uniformly importing individual-factor
operators into the identical-matter completion requires the
crossed metric, including its singular normalization.

There is useful newer literature here. The LS general channel
overlap was originally supported by a finite-level conjecture;
HVX derive the corresponding fusion-channel swapping weight
from an Askey--Wilson transform \cite{HVX}, (82)--(83).
This updates the project's literature map. It does not by itself
identify swapping with a unitary scattering matrix or resolve
our identical-matter moving limit. We leave the older sealed
calculation intact and use the newer result as a subsequent input.

Even after ordered addition, an opposite momentum is not fixed.
The spectral representation of $Q_\Sigma,T_\Sigma$ can be written
as $L^2(\mathbb R_+,dP)\otimes\mathcal K$, with
$Q_\Sigma=P$ and $T_\Sigma=i(P\partial_P+1/2)$.
Here $\mathcal K$ includes the relative coordinate and, if retained,
the common multiplicity. A positive self-adjoint operator
strongly commuting with $P$ and obeying the opposite dilation
covariance is of the form
\begin{equation}
 P^- = \frac{A}{\mu P},\qquad P^+=\mu P,\qquad A\ge0
 \text{ on }\mathcal K.\label{mass}
\end{equation}
Indeed strong commutation makes it a measurable operator field
in $P$; dilation covariance makes $P P^-(P)$ constant, up to
null sets. Equivalently the logarithmic variable makes this a
translation-invariant field. Appropriate self-adjoint domains
are understood through the joint spectral calculus.
Thus the internal operator $A=P^+P^-$ remains free kinematic data.
For two freely assigned masses its usual choice is
$A=m_1^2/x+m_2^2/(1-x)$, $x=p_1/(p_1+p_2)$;
the chord boost and positive $Q$ alone do not select it.

The other LS limit $S=E_\Sigma+P_\Sigma$ obeys
$[Q_\Sigma,S]=-2iT_\Sigma$, and hence is not a commuting
opposite Poincare momentum. This observation does not require
global de Sitter translations to commute; it concerns the
specific flat-space particle interpretation sought here.
Selected-state ordering, four-bilocal dynamics and finite-$G$
effects must still supply additional physical information.

\newpage
\section{Verification and next target}
Independent rational Gram recursion and the triangular images
agree in 1,484 entries, with 280 outer lowering actions checked.
Parameters are $q=1/4,4/9$, $(d_1,d_2)=(2,2),(1,2)$ and total
grades zero through four. Monic orthogonal coordinates make
these comparisons rational and avoid numerical Gram degeneracy.

The moving check uses $m=8,32,128,512$, $\eps=m^{-3}$,
$(d_1,d_2)=(2,2)$, and
$(j,k_1,k_2)=(0,0,0),(1,1,0),(2,1,2)$.
Eight words and both $[\B,\PP]-iE$ and $[\Bt,Q]-iQ$ are tested,
with $Q=E-\PP$. Common cutoffs differ by 48, and all axes are
padded beyond the word shifts. The input tail is bounded by a
polynomial-weighted Poisson Chernoff estimate; it does not on
its own certify unbounded-word errors.

\begin{center}
\begin{tabular}{rrr}
\toprule
$m$ & relative $Q$ error & $\|([\Bt,Q]-iQ)\V\ket{2,1,2}\|$\\
\midrule
8 & 0.200467 & 3.222535\\
32 & 0.090199 & 1.181758\\
128 & 0.042464 & 0.414286\\
512 & 0.020876 & 0.173320\\
\bottomrule
\end{tabular}
\end{center}
Here the relative $Q$ error uses
$\|Q_\Sigma\ket{2,1,2}\|$ as denominator. At $m=128$ the
same nontrivial descendant was independently evaluated with
60-decimal sparse outer-ladder composition. This uses the direct
LS definition, whereas the double-precision code uses
\eqref{stable}. Eight quantities were compared, including the
two commutator defects themselves; the largest relative vector
discrepancy was $5.93\times10^{-13}$. Full cutoffs, tail bounds,
cancellation ratios, versions and input/output hashes accompany
the data. Refinement and high-precision agreement are diagnostics,
not interval arithmetic certificates.

The next target is the crossed metric in this explicit moving
frame, followed by its effect on the actual identical-singlet
four-point function. This preserves the singular first relative
channel rather than discarding it as a free-limit null vector.
The opposite momentum and mass require that dynamical step;
repeating the standard tensor decomposition will not determine them.

\begin{thebibliography}{9}
\bibitem{LS} H.~W. Lin and D. Stanford,
\emph{A symmetry algebra in double-scaled SYK},
\href{https://arxiv.org/abs/2307.15725}{arXiv:2307.15725},
Secs.~2.3, 3.3 and Apps.~B.2, F.
\bibitem{HVX} J. van der Heijden, E. Verlinde and J. Xu,
\emph{Quantum Symmetry and Geometry in Double-Scaled SYK},
\href{https://arxiv.org/abs/2511.08743v2}{arXiv:2511.08743v2},
Secs.~3.1--3.4, especially (55), (62), (82)--(83), (92).
\end{thebibliography}
\end{document}
